\section*{Hoofdstuk \ref{ch:b2:magmatic-differentiation} --- Magmatische differentiatie}

\begin{solution}{exo:b2:magmatic-differentiation:1}
(a) De liquidus van An$_{50}$ ligt op \qty{1450}{\celsius}. Het eerste kristal ligt bij die
temperatuur op de solidus, tussen An$_{66.7}$ (1372) en An$_{83.3}$ (1465):
$66.7 + 16.6\times(1450 - 1372)/(1465 - 1372) = 66.7 + 13.9 = \qty{80.6}{\percent}$
anorthiet. (b) Bij \qty{1394}{\celsius} ligt de vloeistof op de liquidus,
An$_{33.3}$; de kristallen op de solidus, $66.7 + 16.6\times22/93 =
\qty{70.6}{\percent}$. Volgens de hefboomregel is de gekristalliseerde fractie $(50 -
33.3)/(70.6 - 33.3) = 16.7/37.3 = 0.448$. (c) De laatste vloeistof verdwijnt wanneer de
kristallen An$_{50}$ bereiken, bij de solidustemperatuur \qty{1287}{\celsius}; de
liquidus ligt daar tussen An$_{11.1}$ (1265) en An$_{20}$ (1334): $11.1 +
8.9\times22/69 = \qty{13.9}{\percent}$ anorthiet.
\end{solution}

\begin{solution}{exo:b2:magmatic-differentiation:2}
$\bar D_{\ce{Sr}} = 0.15\times0.000063 + 0.35\times0.157 + 0.50\times2.7 = 0.00001 +
0.0550 + 1.35 = 1.40$; $\bar D_{\ce{Zr}} = 0.0001 + 0.0683 + 0.0005 = 0.069$;
$\bar D_{\ce{Rb}} = 0.00001 + 0.0012 + 0.0125 = 0.014$. Alleen strontium, vastgehouden door het
plagioklaas, is compatibel ($\bar D > 1$); zirkonium en rubidium zijn incompatibel,
rubidium het sterkst.
\end{solution}

\begin{solution}{exo:b2:magmatic-differentiation:3}
(a) $C_L/C_0 = 0.5^{D-1}$: Sr $0.5^{0.40} = 0.758$ (het daalt); Zr $0.5^{-0.931} =
1.907$; Rb $0.5^{-0.986} = 1.981$, dicht bij de limiet 2. (b) Zr/Rb verandert met
$0.5^{0.069 - 0.014} = 0.5^{0.055} = 0.963$, minder dan \qty{4}{\percent}, terwijl elk
element bijna verdubbelde: zoals het gevolg zegt, behoudt de verhouding van twee \omterm{def:b2:elements-in-the-earth:partition}{incompatibele
elementen} de waarde van het oudermagma, en zegt ze dus iets over de bron, niet
over de kristallisatie.
\end{solution}

\begin{solution}{exo:b2:magmatic-differentiation:4}
(a) \ce{SiO2} $+5.03$, \ce{MgO} $-1.95$, \ce{K2O} $+0.82$ (procent), kali
vermenigvuldigd met $1.67/0.85 = 1.96$. (b) Voor een volmaakt \omterm{def:b2:elements-in-the-earth:partition}{incompatibel element} geldt
$C_L/C_0 = 1/F$, dus $F = 0.85/1.67 = 0.509$: de helft van het \omterm{def:b2:magma-emplacement:magma}{magma} had moeten
kristalliseren. (c) \qty{12.5}{g} per \qty{100}{g} verwijderen laat $F \approx 0.88$ over en
zou kali maar met ongeveer $1/0.88 = 1.14$ verhogen: kristallisatie alleen kan de
stijging niet verklaren. De kali kwam van buiten, uit de granitische xenolieten, rijk aan
kali (\qty{3.6}{\percent}), die het \omterm{def:b2:magma-emplacement:magma}{magma} assimileerde.
\end{solution}

\begin{solution}{exo:b2:magmatic-differentiation:5}
(a) $F = 0.8$, $D = 5$. Fractioneel: $0.8^{4} = 0.410$; evenwicht:
$1/(5 + 0.8\times(1 - 5)) = 1/1.8 = 0.556$. (b) $F = 0.1$, $D = 0.01$. Fractioneel:
$0.1^{-0.99} = 9.77$; evenwicht: $1/(0.01 + 0.1\times0.99) = 1/0.109 = 9.17$.
(c) Voor $F \to 0$ gaat de wet van Rayleigh naar $0$ voor $D > 1$ en naar oneindig voor
$D < 1$, terwijl de evenwichtswet naar $1/D$ gaat (de laatste vloeistof in evenwicht
met kristallen van de totale samenstelling). Alleen \omterm{def:b2:magmatic-differentiation:crystallisation}{fractionele kristallisatie} kan een
vloeistof volledig van een \omterm{def:b2:elements-in-the-earth:partition}{compatibel element} ontdoen.
\end{solution}

\begin{solution}{exo:b2:magmatic-differentiation:6}
(a) Eén formule-eenheid bevat 1.6 \ce{MgO}, 0.4 \ce{FeO} en één \ce{SiO2}: $1.6\times
40.30 = 64.48$, $0.4\times71.84 = 28.74$, 60.08, totaal \qty{153.30}{g}.
Massapercentages \ce{MgO} \qty{42.06}{\percent}, \ce{FeO} \qty{18.75}{\percent},
\ce{SiO2} \qty{39.19}{\percent}. (b) Hefboomregel op \ce{MgO}: $x = (10.0 - 8.0)/(42.06
- 8.0) = 2.0/34.06 = 0.0587$, ongeveer \qty{6}{\percent} olivijn. (c) Silica van de
vloeistof: $(49.0 - 0.0587\times39.19)/(1 - 0.0587) = 46.70/0.9413 =
\qty{49.61}{\percent}$. Het verwijderen van olivijn, armer aan silica dan basalt, verhoogt het
silica van de vloeistof, maar slechts langzaam.
\end{solution}

\begin{solution}{exo:b2:magmatic-differentiation:7}
(a) $F(0) = 0.768$; $F(1000) = 0.1091 - 0.5906 + 0.7678 = 0.286$; $F(2000) = 0.4364
- 1.1813 + 0.7678 = 0.023$. (b) De discriminant is $(5.9064\times10^{-4})^2 -
4\times1.091\times10^{-7}\times0.7678 = 3.4885\times10^{-7} - 3.3507\times10^{-7} =
1.379\times10^{-8}$, met vierkantswortel $1.174\times10^{-4}$; de kleinste wortel is
$H = (5.9064 - 1.174)\times10^{-4}/(2.182\times10^{-7}) = \qty{2169}{m}$, de
Sandwichhorizont. (c) $1 - 0.768 = 0.232$: bijna een kwart van het \omterm{def:b2:magma-emplacement:magma}{magma} was
gekristalliseerd voordat de basis van de dagzomende Onderste Zone ontstond, in de Verborgen Zone
eronder en in de Randseries langs de wanden en het dak.
\end{solution}

\begin{solution}{exo:b2:magmatic-differentiation:8}
Het \omterm{def:b2:magmatic-differentiation:cumulate}{cumulaat} bevat $D\,C_L = D\,C_0F^{D-1}$, dus tussen twee niveaus geldt
$\ln(\mathrm{Ni}_2/\mathrm{Ni}_1) = (D - 1)\ln(F_2/F_1)$. (a) $D - 1 = \ln(54/177)/
\ln(0.29/0.70) = (-1.187)/(-0.881) = 1.35$, dus $D \approx 2.35$ (de regressie over
de 66 \omterm{def:b2:magmatic-differentiation:cumulate}{cumulaten} met $F \ge 0.23$ geeft 2.38). (b) $D - 1 = \ln(6.2/54)/\ln(0.17/0.29) = (-2.164)/(-0.534)
= 4.05$, dus $D \approx 5.05$. (c) De globale $D$ verandert wanneer de kristalliserende
associatie verandert (een nieuw mineraal dat nikkel opneemt, of een dat verdwijnt), en
de coëfficiënten mineraal--smelt zelf stijgen naarmate de vloeistof afkoelt en van
samenstelling verandert, zoals voor olivijn gemeten; gevangen vloeistof in de \omterm{def:b2:magmatic-differentiation:cumulate}{cumulaten} vertroebelt de
schatting.
\end{solution}

\begin{solution}{exo:b2:magmatic-differentiation:9}
(a) $100 + 25.4 - 2.9 - 9.6 = \qty{112.9}{g}$. (b) Grammen \ce{SiO2}: $55.53 +
25.4\times0.7248 - 2.9\times0.3919 - 9.6\times0.5054 = 55.53 + 18.41 - 1.14 - 4.85 =
\qty{67.95}{g}$, dus $67.95/112.9 = \qty{60.19}{\percent}$ tegen 60.16. \ce{MgO}:
$4.66 + 0.20 - 1.22 = \qty{3.64}{g}$, \qty{3.22}{\percent} tegen 3.19. \ce{K2O}:
$0.88 + 0.92 = \qty{1.80}{g}$, \qty{1.60}{\percent} tegen 1.75. (c) Silica en magnesia, twee van de drie
controle-oxiden, passen binnen 0.03; kali, niet als controle gebruikt, ligt 0.15 te laag, bijna een
tiende van zijn waarde: het geassimileerde materiaal had meer kali dan de gemiddelde
xenoliet, of er werd bij voorkeur een kalirijke smelt ervan opgenomen.
\end{solution}

\begin{solution}{exo:b2:magmatic-differentiation:10}
(a) $z = (0.2 + 1.5 - 1)/(0.2 - 1) = 0.7/(-0.8) = -0.875$; $F^{-z} = 0.7^{0.875} =
0.7319$; $C^*/C_0 = 0.2\times0.5/(0.2 - 1 + 1.5) = 0.1/0.7 = 0.1429$, dus $C/C_0 =
0.7319 + 0.1429\times(1 - 0.7319) = 0.7319 + 0.0383 = 0.770$. (b) Rayleigh:
$0.7^{0.5} = 0.837$. Het \omterm{def:b2:magma-emplacement:magma}{AFC-magma} heeft om twee redenen minder strontium: om
\qty{30}{\percent} van zijn massa te verliezen terwijl het $r$ keer wint wat het kristalliseert, moet het
$(1 - 0.7)/(1 - 0.2) = 0.375$ kristalliseren in plaats van 0.30, en het strontium dat het wint,
komt uit een gesteente dat half zo rijk is. (c) $(0.7319\times0.7035 + 0.0383\times0.7200)/0.7702
= (0.5149 + 0.0276)/0.7702 = 0.70432$: de verhouding is met 0.0008 gestegen, terwijl
kristallisatie alleen haar op 0.7035 zou hebben gelaten.
\end{solution}

\begin{solution}{exo:b2:magmatic-differentiation:11}
(a) $1.1\times(1100 - 200) + 250 = 990 + 250 = \qty{1240}{J}$ per gram. (b)
$25.4\times1240 = \qty{31496}{J}$, dicht bij de \qty{31350}{J} van Wilcox. (c)
$31496/370 = \qty{85.1}{g}$ plagioklaas. (d) De massabalans verwijderde maar $2.9 +
9.6 = \qty{12.5}{g}$ kristallen, waarvan de latente warmte (ongeveer \qty{4600}{J}) minder
dan een zesde van de behoefte dekt; $r = 25.4/12.5 = 2.03$, tweemaal zoveel gesteente geassimileerd als
kristallen gevormd, terwijl het warmtebudget een $r$ ruim onder 1 toelaat. Het uitgebarsten \omterm{def:b2:magma-emplacement:magma}{magma}
kan de warmte niet zelf hebben geleverd: een grotere massa \omterm{def:b2:magma-emplacement:magma}{magma} moet in de diepte,
buiten het zicht, zijn gekristalliseerd en afgekoeld.
\end{solution}

\begin{solution}{exo:b2:magmatic-differentiation:12}
(a) Etappe 1, $F$ van 1 tot 0.8 met $D = 5$: factor $0.8^{4} = 0.4096$. Etappe 2,
van 0.8 tot 0.5 met $D = 1.5$: de wet geldt voor de overgebleven vloeistof, dus de
factor is $(0.5/0.8)^{0.5} = 0.7906$. Nikkel: $250\times0.4096\times0.7906 =
\qty{81.0}{ppm}$. (b) $250\times0.5^{4} = \qty{15.6}{ppm}$, vijf keer te laag. (c)
De wet van Rayleigh veronderstelt een constante $D$; wanneer de kristalliserende associatie verandert,
springt $D$, en moet de wet opnieuw beginnen vanaf de samenstelling van de vloeistof die bij de
verandering is bereikt, met $F$ gemeten ten opzichte van die vloeistof.
\end{solution}

\begin{solution}{pb:b2:magmatic-differentiation:1}
\noindent\textbf{Deel I.}

\textbf{1.} $\bar D_{\ce{Ni}} = 0.20\times5.0 + 0.50\times0 + 0.30\times2.0 = 1.60$.

\textbf{2.} $\bar D_{\ce{Zr}} = 0.20\times0.0007 + 0.50\times0.001 + 0.30\times0.195 =
0.00014 + 0.0005 + 0.0585 = 0.059$.

\textbf{3.} Bovenste zone: $\bar D_{\ce{Ni}} = 0.45\times2.0 = 0.90$; $\bar D_{\ce{Zr}} =
0.55\times0.001 + 0.45\times0.195 = 0.088$.

\textbf{4.} Nikkel is compatibel in de onderste zone (1.60) en wordt licht
incompatibel (0.90) wanneer olivijn verdwijnt; zirkonium is in beide incompatibel.

\textbf{5.} Clinopyroxeen: het levert $0.0585$ van de $0.059$ van de onderste zone, bijna
alles, omdat het het enige mineraal is met een coëfficiënt ver van nul.

\textbf{6.} $(C_{\ce{Ni}}/C_{\ce{Zr}})$ verandert met $F^{D_{\ce{Ni}} - D_{\ce{Zr}}}$;
per halvering $0.5^{1.60 - 0.059} = 0.5^{1.541} = 0.344$.

\noindent\textbf{Deel II.}

\textbf{7.} $1.60\times150 = \qty{240}{ppm}$.

\textbf{8.} De vloeistof had $190/1.60 = \qty{118.75}{ppm}$; $F = (118.75/150)^{1/0.6} =
0.7917^{1.667} = 0.68$.

\textbf{9.} $80\times0.68^{-0.941} = 80\times1.438 = \qty{115}{ppm}$.

\textbf{10.} Fractioneel: $150\times0.68^{0.6} = 150\times0.793 = \qty{119.0}{ppm}$;
evenwicht: $150/(1.6 + 0.68\times(1 - 1.6)) = 150\times0.839 = \qty{125.8}{ppm}$.
Het verschil is maar \qty{6}{\percent} nadat een derde van het \omterm{def:b2:magma-emplacement:magma}{magma} is
gekristalliseerd: voor een matig \omterm{def:b2:elements-in-the-earth:partition}{compatibel element} blijven de twee wetten in het begin dicht bij elkaar.

\textbf{11.} $150\times0.55^{0.6} = \qty{104.8}{ppm}$.

\textbf{12.} Net eronder: $1.60\times104.8 = \qty{167.7}{ppm}$; net erboven:
$0.90\times104.8 = \qty{94.3}{ppm}$. Het nikkelprofiel van de \omterm{def:b2:magmatic-differentiation:cumulate}{cumulaten} daalt
abrupt met \qty{44}{\percent} op het niveau waar olivijn uit de
lagen verdwijnt: een trap, geen vloeiende kromme.

\textbf{13.} Vloeistof: $104.8\times(0.40/0.55)^{0.90 - 1} = 104.8\times1.032 =
\qty{108.2}{ppm}$; \omterm{def:b2:magmatic-differentiation:cumulate}{cumulaat} $0.90\times108.2 = \qty{97.4}{ppm}$. Met $\bar D < 1$
stijgt nikkel nu langzaam in de vloeistof en in de \omterm{def:b2:magmatic-differentiation:cumulate}{cumulaten}: het profiel klimt na zijn
trap zachtjes.

\noindent\textbf{Deel III.}

\textbf{14.} $x_{\mathrm{ol}}(m_{\mathrm{ol},j} - d_j) + x_{\mathrm{pl}}
(m_{\mathrm{pl},j} - d_j) + x_{\mathrm{cpx}}(m_{\mathrm{cpx},j} - d_j) = 100\,(p_j -
d_j)$.

\textbf{15.} \ce{MgO}: $37.16\,x_{\mathrm{ol}} - 4.9\,x_{\mathrm{pl}} +
11.1\,x_{\mathrm{cpx}} = 410$; \ce{CaO}: $-7.3\,x_{\mathrm{ol}} +
7.06\,x_{\mathrm{pl}} + 13.2\,x_{\mathrm{cpx}} = 420$; \ce{Al2O3}:
$-14.1\,x_{\mathrm{ol}} + 17.6\,x_{\mathrm{pl}} - 11.1\,x_{\mathrm{cpx}} = 190$.

\textbf{16.} Oplossen (door eliminatie of met de regel van Cramer): $x_{\mathrm{ol}} = 9.03$,
$x_{\mathrm{pl}} = 30.84$, $x_{\mathrm{cpx}} = 20.32$ gram per \qty{100}{g}.

\textbf{17.} Kristallen $9.03 + 30.84 + 20.32 = \qty{60.19}{g}$; vloeistof
\qty{39.81}{g}: $F = 0.398$.

\textbf{18.} $(4950 - 9.03\times39.19 - 30.84\times50.54 - 20.32\times51.5)/39.81 =
(4950 - 353.9 - 1558.7 - 1046.5)/39.81 = 1990.9/39.81 = \qty{50.0}{\percent}$, zoals in
de \omterm{def:b2:magma-emplacement:bodies}{gang}: het oxide dat niet als controle werd gebruikt, wordt gereproduceerd.

\textbf{19.} De verwijderde kristallen zijn $9.03/60.19 = \qty{15}{\percent}$ olivijn,
\qty{51}{\percent} plagioklaas en \qty{34}{\percent} clinopyroxeen, tussen de
modi van de twee zones in. Ze komen overeen met de onderste associatie die kristalliseert van $F = 1$ tot
0.55 ($0.45\times20 = 9.0$, $0.45\times50 = 22.5$, $0.45\times30 = 13.5$), gevolgd door
de bovenste van 0.55 tot 0.40 ($0.15\times55 = 8.25$, $0.15\times45 = 6.75$): 9.0,
30.75 en 20.25 gram.

\textbf{20.} $0.30/0.75 = 0.40$.

\noindent\textbf{Deel IV.}

\textbf{21.} $0.55^{0.059 - 1} = 0.55^{-0.941} = 1.755$; $80\times1.755 =
\qty{140.4}{ppm}$.

\textbf{22.} $188/140.4 = 1.339 = (F/0.55)^{0.088 - 1}$, dus $F = 0.55\times1.339^{-1/
0.912} = 0.55\times0.726 = 0.399$.

\textbf{23.} $(188/80)^{1/(0.07 - 1)} = 2.35^{-1.075} = 0.399$. Voor \omterm{def:b2:elements-in-the-earth:partition}{incompatibele
elementen} geldt $C_L/C_0 \approx 1/F$, wat de kleine $D$ ook is: het antwoord hangt nauwelijks af van
de details van $D$, wat \omterm{def:b2:elements-in-the-earth:partition}{incompatibele elementen} tot de beste meters van $F$ maakt.

\textbf{24.} Plagioklaas, dat geen ijzer bevat, vormt de helft van de kristallen, en
pyroxeen en olivijn samen bevatten minder ijzer dan magnesia: de vloeistof
wordt rijker aan ijzer en armer aan magnesia zonder silica te winnen. Deze
ijzerverrijking bij bijna constant silica is het kenmerk van de tholeïtische reeks van het
volgende hoofdstuk.

\textbf{25.} Zuivere cumuluskristallen zouden $0.088\times188 = \qty{16.5}{ppm}$ bevatten. Met een
massafractie $t$ gevangen vloeistof geldt $16.5(1 - t) + 188\,t = 47$, dus $t = (47 -
16.5)/(188 - 16.5) = 30.5/171.5 = 0.18$: ongeveer \qty{18}{\percent} gevangen vloeistof.

\textbf{26.} De drie schattingen, hoofdelementen (0.398), kali (0.40) en
zirkonium (0.399), stemmen overeen: toen de bovenste laag kristalliseerde, was
\textbf{\qty[text-series-to-math=true]{40}{\percent}} van het oorspronkelijke \omterm{def:b2:magma-emplacement:magma}{magma}
nog vloeibaar.
\end{solution}
